\documentclass{amsart}
% For dealing with references we use the comment environment
\usepackage{verbatim}
\newenvironment{reference}{\comment}{\endcomment}
%\newenvironment{reference}{}{}
\newenvironment{slogan}{\comment}{\endcomment}
\newenvironment{history}{\comment}{\endcomment}

% For commutative diagrams we use Xy-pic
\usepackage[all]{xy}

% We use 2cell for 2-commutative diagrams.
\xyoption{2cell}
\UseAllTwocells

% We use multicol for the list of chapters between chapters
\usepackage{multicol}

% This is generally recommended for better output
\usepackage{lmodern}
\usepackage[T1]{fontenc}

% For cross-file-references

% Package for hypertext links:
\usepackage{hyperref}

% For any local file, say "hello.tex" you want to link to please

% Theorem environments.
%
\theoremstyle{plain}
\newtheorem{theorem}[subsection]{Theorem}
\newtheorem{proposition}[subsection]{Proposition}
\newtheorem{lemma}[subsection]{Lemma}

\theoremstyle{definition}
\newtheorem{definition}[subsection]{Definition}
\newtheorem{example}[subsection]{Example}
\newtheorem{exercise}[subsection]{Exercise}
\newtheorem{situation}[subsection]{Situation}

\theoremstyle{remark}
\newtheorem{remark}[subsection]{Remark}
\newtheorem{remarks}[subsection]{Remarks}

\numberwithin{equation}{subsection}

% Macros
%
\def\lim{\mathop{\mathrm{lim}}\nolimits}
\def\colim{\mathop{\mathrm{colim}}\nolimits}
\def\Spec{\mathop{\mathrm{Spec}}}
\def\Hom{\mathop{\mathrm{Hom}}\nolimits}
\def\Ext{\mathop{\mathrm{Ext}}\nolimits}
\def\SheafHom{\mathop{\mathcal{H}\!\mathit{om}}\nolimits}
\def\SheafExt{\mathop{\mathcal{E}\!\mathit{xt}}\nolimits}
\def\Sch{\mathit{Sch}}
\def\Mor{\mathop{\mathrm{Mor}}\nolimits}
\def\Ob{\mathop{\mathrm{Ob}}\nolimits}
\def\Sh{\mathop{\mathit{Sh}}\nolimits}
\def\NL{\mathop{N\!L}\nolimits}
\def\CH{\mathop{\mathrm{CH}}\nolimits}
\def\proetale{{pro\text{-}\acute{e}tale}}
\def\etale{{\acute{e}tale}}
\def\QCoh{\mathit{QCoh}}
\def\Ker{\mathop{\mathrm{Ker}}}
\def\Im{\mathop{\mathrm{Im}}}
\def\Coker{\mathop{\mathrm{Coker}}}
\def\Coim{\mathop{\mathrm{Coim}}}

% Boxtimes
%
\DeclareMathSymbol{\boxtimes}{\mathbin}{AMSa}{"02}

%
% Macros for moduli stacks/spaces
%
\def\QCohstack{\mathcal{QC}\!\mathit{oh}}
\def\Cohstack{\mathcal{C}\!\mathit{oh}}
\def\Spacesstack{\mathcal{S}\!\mathit{paces}}
\def\Quotfunctor{\mathrm{Quot}}
\def\Hilbfunctor{\mathrm{Hilb}}
\def\Curvesstack{\mathcal{C}\!\mathit{urves}}
\def\Polarizedstack{\mathcal{P}\!\mathit{olarized}}
\def\Complexesstack{\mathcal{C}\!\mathit{omplexes}}
% \Pic is the operator that assigns to X its picard group, usage \Pic(X)
% \Picardstack_{X/B} denotes the Picard stack of X over B
% \Picardfunctor_{X/B} denotes the Picard functor of X over B
\def\Pic{\mathop{\mathrm{Pic}}\nolimits}
\def\Picardstack{\mathcal{P}\!\mathit{ic}}
\def\Picardfunctor{\mathrm{Pic}}
\def\Deformationcategory{\mathcal{D}\!\mathit{ef}}

\setlength{\emergencystretch}{3em}
\begin{document}
\title[Stacks Project, Tag 0E30]{262 --- Composition of rigidified relative dualizing complexes for flat morphisms}
\maketitle
\noindent Copyright (C) 2005--2025 Johan de Jong. Stacks Project authors. Source: \href{https://stacks.math.columbia.edu/tag/0E30}{Tag 0E30}.

\noindent Editorial difficulty note (not author text). Difficulty H2: The proof constructs the required rigidification for the tensor-product candidate through two successive diagonals. It combines exact-support derived Hom, tensor compatibility and flat base change, then repeats the construction on diagonal neighbourhoods in the nonseparated case. The rigidifying map is an actual extra construction beyond algebraic tensor transitivity..

\noindent Editorial provenance note (not author text). History: excerpt from revision \texttt{a0444\allowbreak 6e57e\allowbreak c1fbc\allowbreak 252a8\allowbreak 71afc\allowbreak ec775\allowbreak 2fb28\allowbreak 07b14}, duality.tex, lines 7405--7494. Reference presentation was adapted; original-author context and core support blocks were retained or added. The mathematical wording is unchanged. Licensed under GNU FDL 1.2 or later; no invariant sections or cover texts. See ../licenses/COPYING. Context blocks reproduce author definitions and supporting results; they are not additional counted problems.

\begin{definition}
\label{definition-relative-dualizing-complex}
Let $X \to S$ be a morphism of schemes which is flat and
locally of finite presentation. Let $W \subset X \times_S X$
be any open such that the diagonal $\Delta_{X/S} : X \to X \times_S X$
factors through a closed immersion $\Delta : X \to W$.
A {\it relative dualizing complex} is a
pair $(K, \xi)$ consisting of an object $K \in D(\mathcal{O}_X)$
and a map
$$
\xi : \Delta_*\mathcal{O}_X \longrightarrow L\text{pr}_1^*K|_W
$$
in $D(\mathcal{O}_W)$ such that
\begin{enumerate}
\item $K$ is $S$-perfect (Derived Categories of Schemes, Definition
\ref{perfect-definition-relatively-perfect}), and
\item $\xi$ defines an isomorphism of $\Delta_*\mathcal{O}_X$
with
$R\SheafHom_{\mathcal{O}_W}(
\Delta_*\mathcal{O}_X, L\text{pr}_1^*K|_W)$.
\end{enumerate}
\end{definition}

\begin{definition}
\label{perfect-definition-relatively-perfect}
Let $f : X \to S$ be a morphism of schemes which is flat and
locally of finite presentation. An object $E$ of $D(\mathcal{O}_X)$ is
{\it perfect relative to $S$} or
{\it $S$-perfect} if $E$ is pseudo-coherent
(Cohomology, Definition \ref{cohomology-definition-pseudo-coherent}) and
$E$ locally has finite tor dimension as an object of
$D(f^{-1}\mathcal{O}_S)$
(Cohomology, Definition \ref{cohomology-definition-tor-amplitude}).
\end{definition}

\begin{definition}
\label{cohomology-definition-pseudo-coherent}
Let $(X, \mathcal{O}_X)$ be a ringed space. Let $\mathcal{E}^\bullet$
be a complex of $\mathcal{O}_X$-modules. Let $m \in \mathbf{Z}$.
\begin{enumerate}
\item We say $\mathcal{E}^\bullet$ is {\it $m$-pseudo-coherent}
if there exists an open covering $X = \bigcup U_i$ and for each $i$
a morphism of complexes
$\alpha_i : \mathcal{E}_i^\bullet \to \mathcal{E}^\bullet|_{U_i}$
where $\mathcal{E}_i^\bullet$ is strictly perfect on $U_i$ and
$H^j(\alpha_i)$ is an isomorphism for $j > m$ and $H^m(\alpha_i)$
is surjective.
\item We say $\mathcal{E}^\bullet$ is {\it pseudo-coherent}
if it is $m$-pseudo-coherent for all $m$.
\item We say an object $E$ of $D(\mathcal{O}_X)$ is
{\it $m$-pseudo-coherent} (resp.\ {\it pseudo-coherent})
if and only if it can be represented by a $m$-pseudo-coherent
(resp.\ pseudo-coherent) complex of $\mathcal{O}_X$-modules.
\end{enumerate}
\end{definition}

\begin{definition}
\label{cohomology-definition-tor-amplitude}
Let $(X, \mathcal{O}_X)$ be a ringed space.
Let $E$ be an object of $D(\mathcal{O}_X)$.
Let $a, b \in \mathbf{Z}$ with $a \leq b$.
\begin{enumerate}
\item We say $E$ has {\it tor-amplitude in $[a, b]$}
if $H^i(E \otimes_{\mathcal{O}_X}^\mathbf{L} \mathcal{F}) = 0$
for all $\mathcal{O}_X$-modules $\mathcal{F}$ and all $i \not \in [a, b]$.
\item We say $E$ has {\it finite tor dimension}
if it has tor-amplitude in $[a, b]$ for some $a, b$.
\item We say $E$ {\it locally has finite tor dimension}
if there exists an open covering $X = \bigcup U_i$ such that
$E|_{U_i}$ has finite tor dimension for all $i$.
\end{enumerate}
An $\mathcal{O}_X$-module $\mathcal{F}$ has {\it tor dimension $\leq d$}
if $\mathcal{F}[0]$ viewed as an object of $D(\mathcal{O}_X)$ has
tor-amplitude in $[-d, 0]$.
\end{definition}

\begin{lemma}
\label{lemma-relative-dualizing-composition}
Let $f : Y \to X$ and $X \to S$ be morphisms of schemes
which are flat and of finite presentation.
Let $(K, \xi)$ and $(M, \eta)$
be a relative dualizing complex for $X \to S$ and $Y \to X$.
Set $E = M \otimes_{\mathcal{O}_Y}^\mathbf{L} Lf^*K$.
Then $(E, \zeta)$ is a relative dualizing complex for $Y \to S$ for
a suitable $\zeta$.
\end{lemma}

\begin{proof}
Using Lemma \href{https://stacks.math.columbia.edu/tag/0E2U}{lemma relative dualizing complex algebra (Tag 0E2U)}
and the algebraic version of this lemma (Dualizing Complexes, Lemma
\href{https://stacks.math.columbia.edu/tag/0E2H}{lemma relative dualizing composition (Tag 0E2H)})
we see that $E$
is affine locally the first component of a relative dualizing complex.
In particular we see that $E$
is $S$-perfect since this may be checked affine locally, see
Derived Categories of Schemes, Lemma
\href{https://stacks.math.columbia.edu/tag/0DI2}{lemma affine locally rel perfect (Tag 0DI2)}.

\medskip\noindent
Let us first prove the existence of $\zeta$ in case the
morphisms $X \to S$ and $Y \to X$ are separated so that
$\Delta_{X/S}$, $\Delta_{Y/X}$, and $\Delta_{Y/S}$
are closed immersions. Consider the following diagram
$$
\xymatrix{
& & Y \ar@{=}[r] & Y \ar[d]^f \\
Y \ar[r]_{\Delta_{Y/X}} &
Y \times_X Y \ar[d]_m \ar[r]_\delta \ar[ru]_q &
Y \times_S Y \ar[d]^{f \times f} \ar[ru]_p & X\\
& X \ar[r]^{\Delta_{X/S}} & X \times_S X \ar[ru]_r
}
$$
where $p$, $q$, $r$ are the first projections.
By Lemma \href{https://stacks.math.columbia.edu/tag/0E2I}{lemma sheaf with exact support internal home (Tag 0E2I)}
we have
$$
R\SheafHom_{\mathcal{O}_{Y \times_S Y}}(
\Delta_{Y/S, *}\mathcal{O}_Y, Lp^*E) =
R\delta_*\left(R\SheafHom_{\mathcal{O}_{Y \times_X Y}}(
\Delta_{Y/X, *}\mathcal{O}_Y,
R\SheafHom(\mathcal{O}_{Y \times_X Y}, Lp^*E))\right)
$$
By Lemma \href{https://stacks.math.columbia.edu/tag/0E2N}{lemma sheaf with exact support tensor (Tag 0E2N)} we have
$$
R\SheafHom(\mathcal{O}_{Y \times_X Y}, Lp^*E) =
R\SheafHom(\mathcal{O}_{Y \times_X Y}, L(f \times f)^*Lr^*K)
\otimes_{\mathcal{O}_{Y \times_S Y}}^\mathbf{L} Lq^*M
$$
By Lemma \href{https://stacks.math.columbia.edu/tag/0E2M}{lemma flat bc sheaf with exact support (Tag 0E2M)} we have
$$
R\SheafHom(\mathcal{O}_{Y \times_X Y}, L(f \times f)^*Lr^*K) =
Lm^*R\SheafHom(\mathcal{O}_X, Lr^*K)
$$
The last expression is isomorphic (via $\xi$) to
$Lm^*\mathcal{O}_X = \mathcal{O}_{Y \times_X Y}$.
Hence the expression preceding is isomorphic to
$Lq^*M$. Hence
$$
R\SheafHom_{\mathcal{O}_{Y \times_S Y}}(
\Delta_{Y/S, *}\mathcal{O}_Y, Lp^*E) =
R\delta_*\left(R\SheafHom_{\mathcal{O}_{Y \times_X Y}}(
\Delta_{Y/X, *}\mathcal{O}_Y, Lq^*M)\right)
$$
The material inside the parentheses is isomorphic to
$\Delta_{Y/X, *}*\mathcal{O}_X$ via $\eta$.
This finishes the proof in the separated case.

\medskip\noindent
In the general case we choose an open $W \subset X \times_S X$
such that $\Delta_{X/S}$ factors through a closed immersion
$\Delta : X \to W$ and we choose an open $V \subset Y \times_X Y$
such that $\Delta_{Y/X}$ factors through a closed immersion
$\Delta' : Y \to V$. Finally, choose an open
$W' \subset Y \times_S Y$ whose intersection with $Y \times_X Y$
gives $V$ and which maps into $W$. Then we consider the diagram
$$
\xymatrix{
& & Y \ar@{=}[r] & Y \ar[d]^f \\
Y \ar[r]_{\Delta'} &
V \ar[d]_m \ar[r]_\delta \ar[ru]_q &
W' \ar[d]^{f \times f} \ar[ru]_p & X\\
& X \ar[r]^\Delta & W \ar[ru]_r
}
$$
and we use exactly the same argument as before.
\end{proof}
\end{document}
